\subsection{Component Contribution Analysis}\label{sec:ablation}

\subsubsection{Aggregate Contributions}

The ablation chain follows a remove-in-sequence paradigm and reads subtractively: each row of Table~\ref{T7} strips one mechanism from the complete KDS-EL (0.9574 accuracy, 88 features on average), and the accuracy loss at each step quantifies the removed component's contribution. These four configurations are exactly the methods compared in Section~\ref{sec:results}: the MI+KER variant is KDS-EL without Two-Phase scheduling, the MI variant is KDS-EL without both the KER strategy and Two-Phase scheduling, and the MEL backbone~\cite{mel2024} is the framework stripped of all three mechanisms---the residual that remains once every component has been removed.

Removing MI-guided initialization and adaptive re-initialization inflicts the largest single accuracy loss: -4.58pp, dropping the mean accuracy from 0.9325 to 0.8868 while the feature count expands from 291 to 2,722, roughly an order of magnitude. This dominant role stems from informed starting positions and stagnation-triggered re-initialization: the MI-informed half provides a high-quality search origin, while adaptive re-initialization counters premature convergence.

Removing the Knowledge-guided Elite Refinement (KER) strategy costs a further -1.63pp, dropping the mean accuracy from 0.9488 to 0.9325 and expanding the feature set from 122 to 291. Because the accumulated weight matrix is derived from fitness comparisons already performed during search, KER's contribution is obtained at a bounded refinement cost that does not alter the dominant per-iteration complexity.

Removing the Two-Phase Weight Scheduling strategy from the complete KDS-EL costs -0.86pp, dropping the mean accuracy from 0.9574 to 0.9488 and expanding the feature set from 88 to 122. Runtime also rises from 51.3\,s to 54.2\,s during this removal.

Each component contributes along a distinct dimension of the search process, and its removal damages that dimension. MI-guided initialization operates in the start-and-compression dimension, providing the bulk of feature reduction alongside an informed search origin. The KER strategy occupies the accuracy dimension, refining elite solutions through cross-iteration accumulated knowledge at a bounded refinement cost. The Two-Phase strategy delivers simultaneous gains across all three axes---accuracy, sparsity, and efficiency---by aligning objective weights with the structural phases of the search. The mechanisms are complementary and independent---none subsumes another, and each is verifiable in isolation.

\subsubsection{Per-Dataset Component Contribution}

Representative rows of Table~\ref{T7} reveal two characteristic patterns across the 20 datasets.

On the eight ceiling datasets---where classification accuracy saturates above 0.998---the MI and KER components rapidly approach the accuracy upper bound, after which the ablative configurations differentiate through feature counts alone. On ALL1, KDS-EL attains 0.9998 accuracy with 39 features; the configuration without Two-Phase scheduling holds 55 features at 0.9996; removing KER in addition leaves the MI-only configuration, also at 55 features; and the backbone selects 4248 features at 0.9664. On Adenoma, the backbone selects 2483 features at 0.9410 accuracy, the MI-only configuration compresses to 120, the configuration without Two-Phase scheduling to 33, and the complete KDS-EL to 66---at the identical saturated accuracy of 0.9996. DLBCL, Leukaemia, Lymphoma, and SRBCT follow the same trajectory: accuracy saturates early, and the configurations differentiate through increasingly compact feature subsets with no accuracy penalty.

On medium-difficulty datasets, accuracy degrades in a stair-step progression as mechanisms are removed. On BASEHOCK, the accuracy sequence reads 0.9113 (KDS-EL), 0.8997 (without Two-Phase scheduling), 0.8597 (additionally without KER), and 0.8518 (backbone). On madelon, a comparable trajectory emerges: 0.8822, 0.8659, 0.7679, and 0.7521. The sharpest feature inflation uniformly occurs when MI-guided initialization is removed (the MI-only-to-backbone step), where feature counts expand by roughly an order of magnitude across nearly all datasets. On Gastric1, features rise from 678 to 7,393; on T1D, from 174 to 6,715; on Prostate, from 27 to 5,058. These patterns confirm that the initial investment in MI-informed initialization delivers both the dominant accuracy contribution and the foundational compression on which subsequent mechanisms build.

\subsubsection{Two-Phase Scheduling and In-place Update Verification}

The Two-Phase Weight Scheduling strategy is isolated by comparing the complete KDS-EL against its fixed-weight counterpart (fixed ws=$[0.9,0.1]$ with in-place update) on the same 20-dataset batch. Removing the schedule lowers accuracy from 0.9574 to 0.9488 (-0.86pp), expands average features from 88 to 122, and increases runtime from 51.3\,s to 54.2\,s. Following the Phase-2 weight switch in the full method, accuracy continues to rise while features undergo further compression, confirming that the coarse-then-fine schedule outperforms a fixed-weight regime. On a per-dataset basis (Table~\ref{T7}), KDS-EL is not inferior to the fixed-weight configuration on the overwhelming majority of datasets.
